\rsection{Gedächtnis Teil 2}
\slide{Begrifflichkeiten}{
    \begin{itemize}
        \item \b{Cue (Hinweisreiz):} Ein Reiz, der den Abruf von Informationen aus dem Gedächtnis auslöst oder erleichtert.
        \item \b{Implizites Gedächtnis:} Beinhaltet erinnerte Inhalte, die nicht bewusst/selbstständig abgerufen werden können, aber das Verhalten oder die Leistung bei Aufgaben beeinflussen.
        \item \b{Kontextspezifisches Gedächtnis:} Das Prinzip der transferangemessenen Verarbeitung besagt, dass die Erinnerungsleistung besser ist, wenn der Kontext der Lernphase mit dem der Testphase übereinstimmt (z. B. Wasser/Land).
        \item \b{Distribution of Practice Effect:} Verteiltes Lernen (mit Pausen zwischen den Durchgängen) führt zu besserem Speichererfolg als massiertes Lernen ("alles auf einmal").
        \item \b{Retrieval Practice Effect:} Der aktive Abruf von Wissen (Testen) wirkt wie ein effektiver Lerndurchgang und stärkt die Gedächtnisspur mehr als reines erneutes Lernen.
    \end{itemize}
}
\slide{Methoden der Gedächtnisforschung}{
    \begin{itemize}
        \item \b{Recall-Varianten:}
        \begin{itemize}
            \item Free Recall: Freie Reproduktion ohne vorgegebene Reihenfolge.
            \item Serial Recall: Wiedergabe muss in der ursprünglichen Präsentationsreihenfolge erfolgen.
            \item Cued Recall: Wiedergabe wird durch Hinweisreize (Cues) unterstützt.
        \end{itemize}
        \item \b{Rekognition vs. Recall:} Rekognition (Wiedererkennen) fällt leichter als Recall, da das Item (Target) selbst als starker Hinweisreiz präsent ist und lediglich ein Bekanntheitsgefühl oder eine Identifikation unter Distraktoren erfordert (Recall benötigt selbstständige Konstruktion von Cues).
        \item Indirekte Methoden (Beispiele):
        \begin{itemize}
            \item Wortanfang ergänzen, Wortfragment lösen, Anagramme lösen, ...
        \end{itemize}
        \item Anwendung bei Frontalhirnschädigung: Diese Methoden werden genutzt, da Patienten oft Störungen im expliziten Abruf (Recall) haben, ihr implizites Gedächtnis (und damit die Leistung in indirekten Tests) jedoch intakt bleibt.
    \end{itemize}
}
\slide{Theorien des Vergessens}{
    \begin{itemize}
        \item \b{Ebbinghaus-Vergessenskurve:} Vergessen verläuft logarithmisch; anfangs wird sehr viel Information sehr schnell vergessen, später flacht die Kurve ab (was lange behalten wurde, bleibt stabil).
        \item \b{Ursachen für Vergessen \& Experimente:}
        \begin{enumerate}
            \item Spurenzerfall (Decay): Die Gedächtnisspur verblasst mit der Zeit (Analogie: Spuren in einem Wachstäfelchen verwaschen).
            \item Interferenz: Ähnliche Informationen stören sich gegenseitig beim Abruf.
            \item Experiment: Lernen von Wortpaaren (A-B). Lernt man danach ähnliche Paare (A-D), sinkt die Erinnerungsleistung für die ersten Paare (retroaktive Interferenz). Lernt man völlig unähnliche Paare (C-D), bleibt die Leistung besser. Dies zeigt, dass Ähnlichkeit zu Störungen führt.
        \end{enumerate}
    \end{itemize}
}
\slide{Speicher-Erfolg: Testing-Effekt}{
    \begin{itemize}
        \item Erklärung: Ein geglückter Abruf (Test) vertieft die Speicherung und ist langfristig effizienter als erneutes Einprägen (Lernen).
        \item Experiment (Roediger \& Karpicke):
        \begin{itemize}
            \item Ablauf: Vergleich von reinem Lernen (SSSS) mit Lernen plus Testen (STTT).
            \item Ergebnisse: Nach kurzer Zeit (5 min) ist reines Lernen (SSSS) besser. Langfristig (nach 1 Woche) ist die Test-Bedingung (STTT) jedoch deutlich überlegen, da der Abruf dem Vergessen entgegenwirkt.
        \end{itemize}
    \end{itemize}
}